\documentclass[11pt,a4paper]{article}
\usepackage[margin=2.2cm]{geometry}
\usepackage[T1]{fontenc}
\usepackage[utf8]{inputenc}
\usepackage{lmodern}
\usepackage{array}
\usepackage{tabularx}
\usepackage{booktabs}
\usepackage{enumitem}
\usepackage{titlesec}
\usepackage{xcolor}
\usepackage{fancyhdr}

\definecolor{heading}{HTML}{2F5D50}

\titleformat{\section}{\Large\bfseries\color{heading}}{}{0pt}{}
\setlength{\parindent}{0pt}
\setlength{\parskip}{0.75em}
\renewcommand{\arraystretch}{1.25}

\begin{document}

\pagestyle{empty}
\begin{center}
    {\Huge\bfseries \textit{West}}\\[0.4em]
    {\Large First Quarter Literature Circle Support}\\[0.8em]
    {\large Focus: Freedom}
\end{center}

\vspace{1em}

\textbf{Text:} Carys Davies, \textit{West}\\
\textbf{Section:} approximately the first quarter of the provided text file, ending around Old Woman From A Distance agreeing to guide Bellman.\\
\textbf{Note:} Page numbers are not available in the \texttt{.txt} file, so references use line numbers from \texttt{Davies\_-Carys-West\_-Scribner\_-libgen.li.txt}.

\newpage

\section{Bard -- Summary}

In the opening section of \textit{West}, John Cyrus Bellman leaves his home in Pennsylvania to travel west in search of the enormous creatures he believes may still exist beyond the frontier. His daughter, Bess, watches him go. Aunt Julie calls him foolish and predicts that he will not return, but Bess sees him very differently: to her he looks brave, purposeful, and almost heroic as he rides away. From the beginning, the theme of freedom is complicated. Bellman's journey looks like freedom because he chooses adventure and discovery over ordinary domestic life, but his choice also takes freedom away from Bess, who must live under Aunt Julie's rules.

The story explains that Bellman became obsessed after reading a newspaper report about huge ancient bones found in Kentucky. The bones make him imagine living monsters somewhere in the unexplored West. He cannot let the idea go. He studies maps and the journals of Lewis and Clark, prepares trade goods, buys a new hat, and asks Julie to take care of Bess. Julie argues that a good father would not abandon his child for such a dangerous dream, but Bellman insists that he has to go and see for himself.

Bellman's journey is difficult but exciting at first. He travels through settlements, crosses the Mississippi, follows the Missouri River, trades with Native people, hunts, fishes, and writes many letters to Bess. His freedom is physical: he moves across rivers, forests, and open country. It is also mental: he is free to follow an idea that no one else believes in. However, winter shows the cost of this freedom. He becomes wet, hungry, cold, and nearly helpless. He survives by eating very little and sleeping wherever he can. Even so, when spring comes, he continues west.

Back home, Bess waits for letters that do not arrive. She talks to Sidney Lott about her father until Sidney insults Bellman and says he will never return. Bess defends her father fiercely and chooses to walk alone rather than forgive Sidney. Her freedom is limited because she is a child, but she still protects her own belief in her father.

The section also introduces Old Woman From A Distance, a young Shawnee man. His people have been forced west after an unfair agreement with the government. Their story shows a darker meaning of ``freedom'': while Bellman chooses to go west, Old Woman's people are pushed west. They lose land, safety, and control over their future. When Old Woman agrees to work for Devereux and then guide Bellman, he is also trying to choose a path for himself in a world where many choices have already been taken from him.

\newpage

\section{Paladin -- Discussion Questions}

\begin{enumerate}[label=\textbf{\arabic*.}, leftmargin=*]
    \item \textbf{Level 1 -- Clarification:} What does Bellman leave behind when he goes west, and what does he take with him?
    \item \textbf{Level 2 -- Explanation:} Why does Bellman feel that he has to go even though Julie, Elmer, and other people think the journey is foolish?
    \item \textbf{Level 3 -- Analysis:} How is Bellman's freedom different from his sister Julie's?
    \item \textbf{Level 4 -- Interpretation:} Bellman chooses to go west, while Old Woman From A Distance's people are forced west. What does this contrast say about freedom in the novel?
    \item \textbf{Level 5 -- Perspective:} What is the cost of freedom in modern society and how does it compare with the context of the novel?
  \end{enumerate}

\newpage

\section{Wizard -- Vocabulary in Context}

\small
\begin{tabularx}{\textwidth}{>{\bfseries}p{2.5cm}p{1.2cm}X X}
\toprule
Word & Line & Short context & Contextual definition \\
\midrule
magnitude & 50 & ``a journey of such magnitude'' & very great size or importance \\
proclamation & 67 & ``in a loud voice like a proclamation'' & a formal public announcement \\
scanty & 92 & ``too scanty or deficient'' & not enough; limited \\
deficient & 92 & ``too scanty or deficient'' & lacking what is needed \\
strangulated & 113 & ``voice high and strangulated'' & tight, strained, or choked-sounding \\
lunatic & 120 & ``a lunatic adventure'' & wildly foolish or unreasonable \\
diverge & 164 & ``then diverge from it'' & to turn away from a route or plan \\
commandeer & 177 & ``commandeer your coffee cup'' & take control of something for one's own use \\
talisman & 294 & ``a kind of talisman against danger'' & an object believed to bring luck or protection \\
arduous & 315 & ``slow and arduous journey'' & difficult and tiring \\
colossal & 371 & ``the colossal bones'' & extremely large \\
prophesied & 471 & ``He prophesied that a time would come'' & predicted, often in a serious or warning way \\
\bottomrule
\end{tabularx}
\normalsize

\newpage

\section{Ranger -- Connections}

This part of \textit{West} connects strongly to stories about characters who chase freedom but discover that freedom has a cost. Bellman reminds me of explorers and dreamers in adventure stories who leave home because ordinary life feels too small. Like characters in frontier films or quest stories, he believes that the unknown world can give his life a greater purpose. His journey west seems exciting because he chooses it himself, studies the route, gathers supplies, and imagines discoveries that no one else has made. In that way, his freedom is connected to imagination.

However, the section also shows that one person's freedom can create unfreedom for someone else. Bellman is free to ride away, but Bess is not free to follow him or stop him. She must stay with Aunt Julie, wait for letters, and defend her father against people who mock him. This connects to real life because adults' choices often affect children, even when the children have no power over those choices. A parent may take a job far away, move the family, or make a risky decision, and the child has to live with the consequences.

Old Woman From A Distance gives the strongest contrast. Bellman goes west because he wants to. Old Woman's people go west because they are pressured and cheated. This reminds me of historical accounts of Native American removal, where ``the West'' did not mean adventure but loss. The same direction can mean freedom for one person and exile for another. That makes the theme more complex: freedom is not just about movement. It is about who gets to choose, who is forced, and who pays the price.

\newpage

\section{Passage Picking}

\begin{enumerate}[label=\textbf{\arabic*.}, leftmargin=*]
    \item \textbf{Important Part -- lines 71--77}\\
    \textbf{Passage to read:} the paragraph where Bess ignores Aunt Julie and watches Bellman ride away, followed by the paragraph where she imagines him as brave, romantic, adventurous, and certain to return.\\
    I chose this passage because it shows Bellman through Bess's eyes. To her, he is not just leaving; he is becoming a heroic figure. It connects to freedom because the West represents possibility for Bellman, even though it also means abandonment and waiting for Bess.

    \item \textbf{Important Part -- lines 207--215}\\
    \textbf{Passage to read:} the exchange between Julie and Bellman where Julie tells him he could simply decide not to go, and Bellman answers that he has to go and see whether the creatures exist.\\
    This passage is one of the clearest explanations of Bellman's motivation. He cannot fully justify the journey with logic, but he feels compelled to make it. The passage is useful for discussing whether freedom means following your own calling, even when others think you are wrong.

    \item \textbf{Good Description -- lines 263--271}\\
    \textbf{Passage to read:} the paragraphs describing Bellman repeatedly reading the newspaper cutting, carrying it close to his heart, imagining the giant animals, and gradually becoming unable to stay where he is.\\
    This passage shows that Bellman's idea has changed him. His home no longer feels like enough. It is a strong freedom passage because it presents freedom as an inner pressure: he feels trapped unless he goes.

    \item \textbf{Surprising / Important Part -- lines 391--427}\\
    \textbf{Passage to read:} the longer section where Bess waits for letters, talks with Sidney Lott about her father's journey, repeats the details of his weapons and trade goods, and finally defends him when Sidney calls him a fool.\\
    I chose this because it shifts attention from Bellman's adventure to Bess's loneliness. It reminds us that Bellman's freedom has emotional consequences for his daughter. It is a good passage for discussing who suffers when someone follows a dream.

    \item \textbf{Important Part -- lines 469--477}\\
    \textbf{Passage to read:} the section where Old Woman From A Distance remembers the men debating whether to accept the unfair agreement and the old man warning that their land and future are being taken from them.\\
    This passage connects freedom to land, survival, and political power. Old Woman's people are not exploring freely; they are being displaced. It deepens the theme by showing that American expansion creates opportunity for some people while taking freedom from others.

    \item \textbf{Important Part -- lines 499--505}\\
    \textbf{Passage to read:} the moment when Devereux offers Old Woman From A Distance profit for helping Bellman, and Old Woman thinks about the old warnings, his sister, what has been lost, and the possible gain before saying yes.\\
    Old Woman's decision to guide Bellman is important because it shows him trying to make a choice inside a limited situation. He may not have full freedom, but he still looks for profit, movement, and a possible future. This makes him more than just a helper in Bellman's story.
\end{enumerate}

\end{document}
